% !TeX program = uplatex
\documentclass[a4paper,12pt]{jlreq}

% ===== Packages =====
\usepackage{amsmath,amssymb,mathtools,bm}
\usepackage{siunitx}
\sisetup{detect-all,per-mode=symbol}
\usepackage{physics}
\usepackage{mhchem}
\usepackage{booktabs}
\usepackage{graphicx}
\usepackage{hyperref}
\usepackage{ascmac}
\usepackage{xcolor}


\hypersetup{
  colorlinks=true,
  linkcolor=black,
  urlcolor=blue,
  citecolor=black
}
\usepackage{enumitem}
\usepackage{geometry}
\usepackage{fancybox}
\geometry{margin=25mm}

% ===== Title =====
\title{Mathematics Notes}
\author{}
\date{\today}

\begin{document}
\maketitle

% ===== Main Content =====
\section*{マクローリン展開とテイラー展開}

\section*{重積分におけるヤコビアンを利用した面積要素の変換 \hfill }
\subsection*{2重積分の場合}

\begin{itembox}[l]{次の微小体積要素の変換を説明せよ}
\begin{align}
\sum_k f(x_k, y_k) A_k 
\;\doteq\;
\sum_k f\bigl(x(u_k, v_k),\, y(u_k, v_k)\bigr)
\left|
\det
\begin{pmatrix}
\dfrac{\partial x}{\partial u} & \dfrac{\partial x}{\partial v} \\[6pt]
\dfrac{\partial y}{\partial u} & \dfrac{\partial y}{\partial v}
\end{pmatrix}
\right|
\Delta u\, \Delta v
\end{align}
\end{itembox}

\subsection*{解答}
一般的な変数変換を考える. いま, 1対1でなめらかな変数変換

\begin{equation}
\begin{cases}
x = x(u, v), \\[4pt]
y = y(u, v)
\end{cases}
\end{equation}
が与えられ, この変換によって, ある領域$D$が長方形領域$D'$になるとする. 

\textcolor{red}{ここに図を挿入} \par

領域$D$の分割によって得られる小領域を$D_k$, その面積を$A_k$とし, 対応する領域$D'$の小領域を$D'_k$とする. 小領域$D'_k$の面積は$\Delta u \Delta v$である. \par

ここで, $\Delta u$や$\Delta v$を用いて$D_k$の面積$A_k$を求めたい. そのためにまず, 3点$(x(u,v),y(u,v)), (x(u + \Delta u, v), y(u + \Delta u, v)), (x(u,v + \Delta v),y(u, v + \Delta v))$をそれぞれP,Q,Rとし, 列ベクトル表示を用いて

\begin{align}
  \overrightarrow{\text{PQ}} = 
  \begin{pmatrix}
    x(u + \Delta u, v) - x(u,v) \\
    y(u + \Delta u, v) - y(u,v)
  \end{pmatrix},
\end{align}

\begin{align}
  \overrightarrow{\text{PR}} =
  \begin{pmatrix}
    x(u, v + \Delta v) - x(u,v) \\
    y(u, v + \Delta v) - y(u,v)
  \end{pmatrix}
\end{align}
によってできる平行四辺形の面積で$A_k$を近似する. \par

2つの列ベクトル$\overrightarrow{\text{PQ}}, \overrightarrow{\text{PR}}$によってできる平行四辺形の面積は$\text{det} (\overrightarrow{\text{PQ}}, \overrightarrow{\text{PR}})$の絶対値であるから

\begin{align*}
A_k &\approx 
\det\begin{pmatrix}
x(u+\Delta u, v) - x(u,v) & x(u, v+\Delta v) - x(u,v) \\
y(u+\Delta u, v) - y(u,v) & y(u, v+\Delta v) - y(u,v)
\end{pmatrix} \\[1em]
&= 
\det\begin{pmatrix}
\dfrac{x(u+\Delta u, v) - x(u,v)}{\Delta u} \, \Delta u & 
\dfrac{x(u, v+\Delta v) - x(u,v)}{\Delta v} \, \Delta v \\[1em]
\dfrac{y(u+\Delta u, v) - y(u,v)}{\Delta u} \, \Delta u &
\dfrac{y(u, v+\Delta v) - y(u,v)}{\Delta v} \, \Delta v
\end{pmatrix}
\end{align*}











\section*{熱力学における変文原理}


\section*{フーリエ変換}


\section*{微分方程式}


\section*{以下は洋書からの抜粋集}










\end{document}
